\section{Half-conformal flatness and the metric classification}
\label{sec:classification}

The coupled estimates have been proved for the standing strictly
positive geometry, including its zero sets. Their strictness now rules
out two nonzero Weyl blocks. We then isolate the remaining harmonic
alternative, exclude it by a conditional lower volume estimate, and
identify the parallel curvature tensor. The orientation-cover argument
comes only after the oriented conclusion has been established.

\begin{proposition}[Half-conformal flatness]\label{prop:one-block-vanishes}
At least one of $T,U$ vanishes identically; equivalently, the metric is
half-conformally flat.
\end{proposition}
\begin{proof}
Suppose otherwise. Proposition~\ref{prop:balanced-moments} gives $\mathbf E v^2=\mathbf E b^2$, so Proposition~\ref{prop:coupling-estimates} applies. Its integrable total expression $\mathcal I$ has nonnegative integral and satisfies
\[
 \mathcal I\le-cj+2\Theta r-\Theta^2/c
 =-\frac{(\Theta-cr)^2}{c}-c(j-r^2),\qquad c\ge\frac1{10},
\]
where
\[
 r=R_v|p|^2+R_b|q|^2,\qquad
 j=(|p|^2+|q|^2)(|p|^2/v+|q|^2/b).
\]
The zero-set conventions established with \eqref{eq:3} apply to the quotients. For $v>0$,
\[
 1+6v+v^2-4\sqrt v(1+v)=(\sqrt v-1)^4,
\]
so $R_v\le1/\sqrt v$, with equality only at $v=1$. Cauchy--Schwarz therefore gives $r^2\le j$.

This inequality is strict almost everywhere on the set where $p$ or $q$ is nonzero. Indeed a Lipschitz function has zero gradient almost everywhere on each of its level sets: at a differentiability and density point of a level set, a nonzero differential would contradict density. Thus, outside a null set, an active gradient $p$ has $v>0$ and $v\ne1$, and similarly for $q$. At such a point
\[
 r< |p|^2/\sqrt v+|q|^2/\sqrt b\le\sqrt j.
\]
Terms with vanishing gradients are read as zero. It follows that $\mathcal I\le0$ almost everywhere, and that $\mathcal I<0$ almost everywhere where either gradient is nonzero. Since $\mathbf E\mathcal I\ge0$, both gradients vanish almost everywhere. Connectedness and continuity make $v,b$ constant. Both constants are positive, so the separate gaps \eqref{eq:5} imply $v+b\ge2$, contradicting $v+b\le x+y<2$.
\end{proof}

Once one Weyl block vanishes, $(M,g)$ belongs to the classically
classified compact half-conformally flat Einstein class with positive
scalar curvature~\cite{Hitchin1981,FriedrichKurke1982,Obata1962}.
We retain the lower volume and parallel curvature argument below as an
independent proof of this final step under our standing hypothesis of
strictly positive sectional curvature.

\subsection{The remaining harmonic alternatives}

After reversing orientation if necessary, one-block vanishing leaves
$U=0$. If $T$ also vanishes, the sectional curvature is constantly one.
For a nonzero surviving block, the characteristic identities reduce the
possibilities to two integer values.

\begin{lemma}[The one-block harmonic alternatives]
\label{lem:one-block-alternatives}
Suppose $U=0$ and $T\not\equiv0$. Then
\[
 \frac{V_M}{\pi^2}=\frac{8-2n_+}{3},\qquad
 s_0=\mathbf E v^2=\frac{3n_+}{4-n_+},\qquad n_+=1\text{ or }2.
\]
In the second case, $s_0=3$ and $V_M=4\pi^2/3$.
\end{lemma}
\begin{proof}
Proposition~\ref{prop:characteristic-identities} gives $n_-=0$ and
$\chi=2+n_+$, $\tau=n_+$. Solving its two identities yields the displayed
volume and moment formulas. Positivity of volume gives $n_+<4$; the
nonzero-block gap gives $s_0\ge1$, while $v<2$ gives $s_0<4$.
The finite nonnegative integer $n_+$ is therefore $1$ or $2$.
Substitution gives the stated values in the second case.
\end{proof}

\subsection{Excluding the exceptional harmonic count}

The value $n_+=2$ would force exactly $s_0=3$ and $V_M=4\pi^2/3$.
We now prove a strict lower bound under that same global moment
hypothesis. The statement does not require one-block vanishing; that
conclusion is used only to produce the exceptional value to which the
bound will be applied.

\begin{proposition}[Conditional lower volume bound]
\label{prop:conditional-lower-volume}
If $s_0=\mathbf E(v^2+b^2)=3$, then $V_M>4\pi^2/3$.
\end{proposition}

Recall from Section~\ref{sec:volume} that $\mathcal E$ is expectation over
the unit tangent bundle, and that $\mathcal R$ is the normal Jacobi
curvature matrix along its geodesics. With
$\mathcal J''+\mathcal R\mathcal J=0$, $\mathcal J(0)=0$ and
$\mathcal J'(0)=\mathrm{Id}$, put $D(t)=\det\mathcal J(t)$.
We use the topology statement in Lemma~\ref{lem:volume-topology}, the
geodesic averages in Lemma~\ref{lem:geodesic-averages}, and the discrete
determinant in Lemma~\ref{lem:discrete-jacobi}.

\begin{lemma}[Lower polar bound]\label{lem:lower-polar}
One has
\[
 V_M\ge2\pi^2\int_0^{\pi/\sqrt3}\mathcal E D(t)\,dt.
\]
\end{lemma}
\begin{proof}
The open ball of radius $\pi/\sqrt3$ lies in the injectivity domain at
every center by Lemma~\ref{lem:volume-topology}. Its volume is the
corresponding integral of $D$ and is at most $V_M$. Averaging over centers
gives the bound. Increasing open radii justify the endpoint; it does not
change the integral.
\end{proof}

\begin{lemma}[Lower radial comparison]\label{lem:lower-radial}
If $s_0=3$ and $0<t<\pi/\sqrt3$, then, with $\theta=\sqrt3t$,
\[
 \mathcal E D(t)\ge
 t^3\left(\frac{\sin\theta}{\theta}\right)^{5/6}
 e^{-\theta^2/36}.
\]
Moreover,
\[
 -\log\left(\frac{\sin\theta}{\theta}\right)-\frac{\theta^2}{6}
 \ge\frac12\left(\frac\theta\pi\right)^4
       +\frac13\left(\frac\theta\pi\right)^6.
\]
\end{lemma}
\begin{proof}
For each integer $N\ge2$, put $a=t/N$. Let $L_0$ be the scalar
Dirichlet tridiagonal matrix on the interior nodes $1,\ldots,N-1$,
with diagonal entries $2$ and adjacent off-diagonal entries $-1$.
Put $G=a^2L_0^{-1}$. Its entries are nonnegative; explicitly,
\[
 G_{ij}=a^2\frac{\min(i,j)(N-\max(i,j))}{N},\qquad
 \operatorname{tr}G=\frac{a^2(N^2-1)}6\longrightarrow\frac{t^2}{6}.
\]
Write $\mathcal R_N$ for the block diagonal matrix with blocks
$\mathcal R(ia)$ and define
$H_N=(G^{1/2}\otimes\mathrm{Id})\mathcal R_N
(G^{1/2}\otimes\mathrm{Id})$.
For fixed $t<\pi/\sqrt3$ and large $N$,
\[
 0\prec H_N\preceq q_t\mathrm{Id},\qquad q_t<1,
 \qquad
 \|H_N\|\le\frac{3a^2}{2-2\cos(\pi/N)},\qquad
 \operatorname{tr}H_N=3\operatorname{tr}G.
\]
In particular,
\[
 -\log\det(\mathrm{Id}-H_N)
 =\sum_{k\ge1}\frac{\operatorname{tr}H_N^k}{k}
 \le\frac{3\operatorname{tr}G}{1-q_t}.
\]
This bound is uniform over all geodesics and large $N$ at the fixed
time.  It justifies taking expectations in the series and gives a
uniform positive lower bound for the determinant ratios, independent
of their increasing matrix dimensions.

By \eqref{eq:8}, $\mathcal E|\mathcal R|^2=3+3s_0/2=15/2$.
For any ordered product of $k\ge2$ curvature matrices,
\[
 |\operatorname{tr}(\mathcal R_1\cdots\mathcal R_k)|
 \le3^{k-2}|\mathcal R_1|\,|\mathcal R_k|.
\]
Cauchy--Schwarz and invariance of geodesic-flow measure bound its
expectation by $(5/6)3^k$.  Neither commutation nor independence of the
matrices is needed.  Expanding the block trace of
$((G\otimes\mathrm{Id})\mathcal R_N)^k$, its scalar coefficients are
products of nonnegative entries of $G$.  Their sum is
$\operatorname{tr}(G^k)$, so
\[
 \mathcal E\operatorname{tr}H_N^k
 \le\frac56\,3^k\operatorname{tr}(G^k)\quad(k\ge2).
\]
The term $k=1$ is exactly $3\operatorname{tr}G$.  Therefore
\[
 \mathcal E\log\det(\mathrm{Id}-H_N)
 \ge\frac56\log\det(\mathrm{Id}-3G)-\frac12\operatorname{tr}G.
\]
Jensen's inequality supplies a lower bound for the expected determinant
by the exponential of this expression.

The flat block determinant is $N^3$, and the normalization is exactly
\[
 \det(\mathrm{Id}-H_N)=\frac{\det\mathcal L_N}{N^3}.
\]
Since $a=t/N$, Lemma~\ref{lem:discrete-jacobi} gives
$\det(\mathrm{Id}-H_N)\to D(t)/t^3$, uniformly over geodesics at this
fixed $t$: on the compact unit tangent bundle the curvature coefficients
and their first derivatives are uniformly bounded.  The scalar version gives
\[
 \det(\mathrm{Id}-3G)\longrightarrow\frac{\sin\theta}{\theta}.
\]
Passing to these limits proves the first assertion.  Equivalently, the
uniform logarithmic bound permits passing averaged logarithms directly;
applying Jensen before the limit already suffices.

Finally, in the scalar positive logarithmic series keep only the
$k=2,3$ terms and, in each, the largest eigenvalue of $G$, which tends
to $t^2/\pi^2$.  Subtracting $3\operatorname{tr}G$ and passing to the
limit yields the second assertion.
\end{proof}

For the exact evaluation, use the new lower-comparison variable
$u=\sqrt3t/\pi$. Lemma~\ref{lem:sine-minorant} gives $a_0=157/50<\pi$
and
\[
 0\le S(\min(u,1-u))\le\sin(\pi u)\qquad(0\le u\le1).
\]
The two strict rational integral bounds below include the prefactor;
the rows use their polynomial representatives up to the shared endpoint.
\begin{center}
\begin{tabular}{@{}llll@{}}
\toprule
Integrand and prefactor & Interval & Degree & Lower bound\\
\midrule
$\dfrac{a_0^3}{27}u^2(u^4/2+u^6/3)S(\min(u,1-u))$
 &$[0,1/2]$&15&$698/10^6$\\
same &$[1/2,1]$&15&$40554/10^6$\\
\bottomrule
\end{tabular}
\end{center}
They follow by the rational polynomial integration formula used in
Section~\ref{sec:volume}. Adding the base term
$2(a_0^2-4)/9=4883/3750$ gives a total greater than
$1343/1000>4/3$.

\begin{proof}[Proof of Proposition~\ref{prop:conditional-lower-volume}]
Write the first bound in
Lemma~\ref{lem:lower-radial} relative to $t^3\sin\theta/\theta$.
Its exponential factor is
\[
 \exp\left[\frac16\left(
 -\log\frac{\sin\theta}{\theta}-\frac{\theta^2}{6}\right)\right]
 \ge1+\frac16\left[\frac12\left(\frac\theta\pi\right)^4
                    +\frac13\left(\frac\theta\pi\right)^6\right].
\]
Integrate Lemma~\ref{lem:lower-polar}, first up to any
$T<\pi/\sqrt3$ and then let $T$ increase to the endpoint.  All integrands
are nonnegative, so monotone convergence suffices; no uniform endpoint
bound for a logarithm is required.  The changes of variable
$\theta=\sqrt3t$, $u=\theta/\pi$ give
\[
 \frac{V_M}{\pi^2}
 \ge\frac29(\pi^2-4)
 +\frac{\pi^3}{27}\int_0^1
       u^2(u^4/2+u^6/3)\sin(\pi u)\,du.
\]
The last two rational integral bounds and their constant term make
this strictly greater than $1343/1000>4/3$, proving the assertion.
\end{proof}

Lemma~\ref{lem:one-block-alternatives} makes the $n_+=2$ case equal to
$V_M=4\pi^2/3$ under $s_0=3$, contrary to
Proposition~\ref{prop:conditional-lower-volume}. Hence only $n_+=1$
remains when exactly one block is nonzero.

\subsection{Parallel curvature and global identification}

\begin{proposition}[The remaining parallel tensors]\label{prop:parallel-models}
Up to reversing orientation, either $T=U=0$, or $U=0$ and $T$ is parallel
with eigenvalues $2,-1,-1$.
\end{proposition}
\begin{proof}
In the nonzero one-block case, the preceding exclusion gives $n_+=1$,
so Lemma~\ref{lem:one-block-alternatives} gives $\mathbf E v^2=1$.
The equality case of Proposition~\ref{prop:weyl-gap} yields
$\nabla T=0$ and $v=\sigma=1$, hence eigenvalues $2,-1,-1$.
The exchanged orientation is identical. If both blocks vanish, the
curvature-operator decomposition gives constant sectional curvature one.
\end{proof}

The model calculation is the Hopf-submersion curvature computation,
following the framework of O'Neill~\cite{ONeill1966}. We include it to
fix the metric scaling precisely.

\begin{lemma}[Curvature of the standard models]\label{lem:model-curvature}
The two tensors in Proposition~\ref{prop:parallel-models} are the curvature tensors of the unit round sphere and of twice the Fubini--Study metric induced by the Hopf quotient of the unit sphere $S^5\subset\mathbb C^3$.
\end{lemma}
\begin{proof}
The first assertion follows immediately from the curvature-operator decomposition. In the second case choose at one point the orthogonal complex structure $I$ corresponding to the eigenline of $T$ with eigenvalue $2$. On an orthonormal pair $X,Y$, the sectional curvature is
\[
 K(X,Y)=\tfrac12\bigl(1+3\langle IX,Y\rangle^2\bigr).
\]
Sectional curvatures determine an algebraic curvature tensor by polarization. Only this choice at a single tangent space is needed; no global complex structure is assumed.

For comparison, let $g_H$ be the quotient metric on $S^5/S^1=\mathbb{CP}^2$ obtained by isometric horizontal lift. At $z\in S^5$, the unit vertical field is $iz$. For horizontal lifts of base vector fields, the horizontal part of the sphere connection is the lifted base connection by the metric formula, while
\[
 (\nabla_XY)^{\mathrm{vert}}=-\langle iX,Y\rangle iz.
\]
This follows by differentiating $\langle Y,iz\rangle=0$. Substitution in the curvature formula yields
\[
 K_{g_H}(X,Y)=1+3\langle iX,Y\rangle^2.
\]
To track the factor three, the vertical bracket is
$[X,Y]^{\mathrm{vert}}=-2\langle iX,Y\rangle iz$ and invariant horizontal lifts satisfy $\nabla_{iz}Y=iY$. The differentiated vertical term and the vertical bracket therefore contribute, respectively, one and two copies of the same squared inner product. Equivalently, the extra horizontal term in the sphere curvature expression is $3\langle iX,Y\rangle iY$, whose pairing with $X$ is $-3\langle iX,Y\rangle^2$ relative to the base expression. Scaling $g_H$ by two gives the required sectional curvatures.

The quotient metric has parallel curvature: a unitary transformation that fixes a complex line and is minus the identity on its orthogonal complement induces an isometry fixing that point with tangent differential $-\operatorname{Id}$. It acts by $-1$ on the covariant five-tensor $\nabla R$, so $\nabla R=0$. Transitivity gives this conclusion at every point. The round sphere likewise has parallel curvature.
\end{proof}

Parallel curvature turns matching tensors at one point into matching
Jacobi equations along paths. This is the continuation principle behind
curvature-based isometry results such as Ambrose's~\cite{Ambrose1956}.
We give the continuation and its inverse explicitly in the present
compact, simply connected setting.

\begin{lemma}[Global continuation of the model isometry]\label{lem:global-isometry}
The local identifications furnished by Lemma~\ref{lem:model-curvature} extend to global isometries.
\end{lemma}
\begin{proof}
Lemma~\ref{lem:volume-topology} shows that $M$ is simply connected. Both model spaces are simply connected as well. For the Hopf quotient this can be seen directly: local sections lift any loop to a path in $S^5$. Join its endpoints in the circle fibre to close the path, contract the resulting loop in $S^5$, and project the contraction. The added fibre path projects to a constant, so the original loop contracts.

Choose a tangent-space isometry matching the full curvature tensor at one point. On sufficiently small normal balls the exponential maps give an isometry: in parallel frames the Jacobi equations have identical coefficients because both curvature tensors are parallel. At a new point, the differential again matches curvature and defines the same construction on a new normal ball. Compactness supplies a common positive radius for these local constructions on the two manifolds.

An isometry germ is determined by its value and differential through exponential maps. Hence it continues uniquely along each path. Subdivide a homotopy square into small pieces lying in normal balls; uniqueness on overlaps shows that continuation is invariant under homotopy. Simple connectivity therefore gives a global local isometry. Continue the inverse germ in the same way on the model. The two compositions agree with the identity near the original points and then everywhere by continuation. They are inverse global isometries.
\end{proof}

\begin{corollary}[The oriented classification]
\label{cor:oriented-classification}
Every connected oriented smooth closed Einstein four-manifold with
strictly positive sectional curvature is, after positive scaling,
isometric to the round $S^4$ or to $\mathbb{CP}^2$ with its
Fubini--Study metric.
\end{corollary}
\begin{proof}
Strict sectional positivity gives a positive Einstein constant, so the
metric can be normalized to $\operatorname{Ric}=3g$.
Proposition~\ref{prop:parallel-models} and
Lemmas~\ref{lem:model-curvature}--\ref{lem:global-isometry} identify
that normalized metric with the unit round sphere or twice the Hopf
quotient metric. Undoing the positive normalization proves the stated
oriented conclusion.
\end{proof}

\subsection{The orientation cover and the third model}

\begin{proof}[Completion of Theorem~\ref{thm:classification}]
Strict sectional positivity makes the Ricci tensor positive definite,
so the Einstein constant is positive.
If $M$ is orientable, choose an orientation and apply
Corollary~\ref{cor:oriented-classification}.

Otherwise let $p:\widetilde M\to M$ be the orientation double cover.
It is connected, closed and oriented; for the standard construction and
connectedness see~\cite[Proposition~3.25]{Hatcher2002}.
The pulled-back metric $\widetilde g=p^*g$ has the same Einstein and
strict sectional-positivity properties. The nontrivial deck map
$\tau$ switches the two local orientations over each point, so it is
a free orientation-reversing involution, and $p\circ\tau=p$ makes it
an isometry of $\widetilde g$. After positive scaling,
Corollary~\ref{cor:oriented-classification} identifies this cover with one of its
two standard models. Conjugating $\tau$ by the model isometry preserves
its freeness and its orientation sign.

The $\mathbb{CP}^2$ case is impossible. Its integral cohomology ring is
$\mathbb Z[h]/(h^3)$ with $|h|=2$, so $h^2$ generates
$H^4(\mathbb{CP}^2;\mathbb Z)$
\cite[Theorem~3.19]{Hatcher2002}. Any self-diffeomorphism sends $h$ to
$\pm h$ and therefore fixes $h^2$. It has degree $+1$, contradicting
the required orientation reversal.

On the unit round $S^4\subset\mathbb R^5$, distance determines ambient
inner products by $\langle x,y\rangle=\cos d(x,y)$. An isometry therefore
preserves these inner products: its values on the five coordinate unit
vectors determine an $A\in O(5)$, and comparison with their images gives
$\tau(x)=Ax$ for every $x\in S^4$. Since $\tau^2=\operatorname{id}$,
we have $A^2=I$. An orthogonal involution has an orthogonal decomposition
into its $+1$ and $-1$ eigenspaces. A nonzero $+1$ eigenvector would give
a fixed point on the sphere, so freeness forces $A=-I$. The scaled base
is consequently the round quotient
$S^4/\{x\sim -x\}=\mathbb{RP}^4$.

The antipodal map has degree $(-1)^5=-1$ on $S^4$, so it reverses
orientation~\cite[Section~2.2, property~(f)]{Hatcher2002}.
Its quotient is a closed connected smooth manifold
and inherits the Einstein and positive-sectional-curvature properties
from the round sphere through the local isometry. Thus this third model
does occur, and it is the only nonorientable case.
\end{proof}
